\section{The Embedding Is the Lever}
\label{sec_disc_embedding}

The most consistent finding across the experiments is that pose-grouping
performance depends more strongly on the quality of the per-keypoint
embedding than on the algorithm used to partition those embeddings into
person identities.
This pattern appears across three separate groups of experiments.
Replacing $k$-means with learned grouping heads does not produce
a consistent improvement,
and the more specialised SA-DMoN objective also fails to improve
the underlying representation.
In contrast, changes made directly to the embedding network produce
gains that are retained, and subsequently strengthened, after transfer
to COCO.

This section considers these results together
rather than as independent negative and positive experiments.
The DEC, slot-attention, and graph-partitioning results show what
happens when additional modelling capacity is placed after or alongside
an embedding that still contains ambiguous person assignments.
The SA-DMoN experiments provide a stronger test of the same idea,
by introducing structural information through the clustering objective,
but again fail to move beyond the performance of the simpler baseline.
PG-GAT changes where that structural information enters the pipeline -
instead of asking the grouping head to recover person structure
from the final embeddings,
it uses joint type,
relative geometry,
and same-type relationships during message passing while those
embeddings are being formed.

Taken together, these experiments suggest that the important
distinction is not simply between simple and learned grouping methods
but between introducing structure \emph{after} the representation has
been formed and using that structure to shape the representation itself.
The remainder of this section examines the evidence for this
interpretation
and relates each PG-GAT modification to the failure modes observed
in the earlier grouping-head experiments.

\subsection{Why every learned head failed on the same embedding}
\label{sec_disc_embedding_heads}

Three learned grouping heads, DEC, slot attention, and graph
partitioning, were attached to the standard GAT backbone of
Section~\ref{sec_meth_baseline} and trained jointly with the embedding
under matched optimiser and schedule settings
(Section~\ref{sec_res_learned_heads}).
None produces a grouping result that improves on the no-depth $k$-means
baseline once evaluated on COCO.
The contrast is clearest for graph partitioning -
it is the strongest learned head on synthetic validation at $0.9714$,
but its own COCO read-out falls to $0.7514$.
Applying $k$-means to the embeddings produced by the same jointly
trained model recovers the score to $0.8366$,
which is still below the standard GAT baseline of $0.8841$.
The three heads fail in different ways,
and this consistency across different grouping mechanisms is what makes
the combined result more informative than any individual negative
experiment.

DEC adds little because its clustering objective has limited ability
to correct errors already present in the embedding.
Its KL-divergence objective compares the soft assignment $q$
with a sharpened target $p$ that is itself constructed from $q$
(Equation~\eqref{eq:dec_p_prelim}).
When $k$-means++ initialisation already produces confident assignments,
the target distribution mainly reinforces those assignments rather than
introducing independent supervisory information about person identity.
This behaviour is particularly limiting for the ambiguous cases that
matter most for pose grouping.
If two same-type keypoints from different people lie close together
in embedding space,
an incorrect initial assignment can be sharpened along with the
correct ones.
DEC therefore has no independent structural cue with which to resolve
the error.
This is consistent with its COCO result of $0.8798$,
which remains close to but below the $0.8841$ baseline.

Slot attention behaves differently.
Its synthetic-validation PGA of $0.9025$ exceeds the standard GAT
baseline of $0.8783$,
showing that a learned competitive assignment mechanism can exploit
the relatively clean synthetic representation.
The improvement does not survive transfer -
its own COCO read-out reaches only $0.7986$,
while $k$-means on the jointly trained embeddings reaches $0.8254$.
One limitation is that the slot competition does not explicitly encode
the type-exclusivity constraint available from the pose representation.
Two detections of the same joint type must belong to different people,
but the slot mechanism is not constructed to enforce this relationship
directly.
On the synthetic distribution the embedding can often resolve the
ambiguity without that prior -
the COCO result suggests that this becomes less reliable once the
representation is exposed to the synthetic-to-real shift.
The experiment does not isolate type-exclusivity as the sole cause
of the transfer failure,
but its absence provides a structural explanation consistent with the
observed behaviour.

Graph partitioning provides the strongest example of a synthetic gain
that does not transfer.
Its $0.9714$ synthetic-validation PGA is well above the $0.8783$
baseline,
yet its own COCO grouping falls to $0.7514$.
The $k$-means read-out on the same jointly trained embeddings recovers
to $0.8366$,
indicating that a substantial part of the failure lies in the
partitioning read-out rather than only in the embedding.
The edge classifier converts pairwise affinities into binary decisions,
which are subsequently thresholded and assembled into connected
components.
This creates a sensitivity that is absent from centroid-based $k$-means:
changes in the affinity distribution around the decision threshold
can alter individual edges,
and a small number of incorrect edges can change the connectivity
of an entire component.
The large synthetic-to-COCO deterioration is therefore consistent with
a partitioning mechanism that is particularly sensitive to distribution
shift.

Taken together, the three experiments provide evidence against the
grouping head being the main limitation of the baseline pipeline.
DEC adds an unsupervised sharpening objective,
slot attention introduces a learned competitive assignment mechanism,
and graph partitioning reformulates grouping as supervised pairwise
classification.
Despite these substantially different approaches,
none produces a COCO result above the $0.8841$ $k$-means baseline.
In several cases, applying $k$-means to the jointly trained embeddings
performs better than the head for which those embeddings were trained.

The common result is therefore more important than the individual
failure mechanisms.
Increasing the complexity of the grouping stage does not reliably
recover information that is absent or ambiguous in the representation
supplied to it.
This observation motivates moving the structural priors earlier
in the pipeline,
so that they influence the construction of the embedding itself
rather than being introduced only when that embedding is partitioned.
The SA-DMoN and PG-GAT experiments provide the next two tests
of this interpretation.

\subsection{Why ten SA-DMoN variants plateaued}
\label{sec_disc_embedding_sadmon}

The SA-DMoN experiments extend the grouping-head investigation in a
different direction.
Rather than learning a direct assignment from the final embeddings,
SA-DMoN introduces a graph-structural objective based on modularity
maximisation over the keypoint $k$NN graph.
The ten configurations reported in Section~\ref{sec_res_sadmon}
all converge to a narrow synthetic-validation band of approximately
$0.77$--$0.79$,
around $0.09$ below the no-depth baseline.
More informative than the best individual result is the fact that
substantial changes to the null model,
adjacency,
and type loss do not move the model outside this band.
This suggests that the limitation lies in the compatibility between
the modularity objective and the graph on which it operates
rather than in a particular SA-DMoN hyperparameter.

Newman modularity (Equation~\eqref{eq:newman_modularity}) measures
whether a proposed partition contains more within-community edge mass
than would be expected under a chosen null model.
This formulation is well suited to graphs in which observed connectivity
carries information about the communities of interest,
and the null model provides a reasonable description of connectivity
in the absence of those communities.
The keypoint $k$NN graph is misaligned with this assumption.
Its edges are constructed explicitly from spatial proximity,
whereas the desired partition is person identity.
Two keypoints can therefore be neighbours because they occupy nearby
image coordinates even when they belong to different people.
This occurs precisely in the crowded and overlapping poses for which
grouping is most difficult.

The consequence is a mismatch between the structure emphasised by the
graph and the structure required by the task.
A modularity objective operating on this graph is encouraged to preserve
patterns of local connectivity,
but those patterns need not correspond to person identity.
In crowded scenes they may instead connect joints belonging to different
people.
The spectral objective is therefore being asked to recover identity
communities from a graph whose edges were not constructed primarily
from identity information.
This provides a structural explanation for why changing the details
of the modularity formulation has little effect on the final grouping
accuracy.

The SA-DMoN variants test this explanation from several directions.
Changing the spatial bandwidth $\sigma$ modifies how strongly proximity
enters the skeleton-aware null model,
but no tested parameterisation produces a useful operating point.
With unconstrained optimisation, $\sigma$ grows to approximately $341$,
making the spatial kernel nearly uniform and effectively weakening
the intended spatial correction.
Constraining $\sigma$ to $[0.05,0.5]$ prevents this divergence,
but the learned value moves to the upper boundary,
again indicating that the optimisation prefers a weaker spatial null.
Fixing $\sigma$ from the per-graph median pairwise distance moves
in the opposite direction and also reduces performance.
The behaviour across these variants is more consistent with a poorly
matched objective than with a single badly chosen bandwidth.

The v2 experiments alter other parts of the formulation but reach the
same conclusion.
Feature-based adjacency breaks the direct dependence between the
message-passing graph and the spatial null
but makes the adjacency itself change as the representation is learned.
The modularity objective must then optimise against a graph that is
moving with the embedding.
Reducing the spectral weight partially recovers the loss in performance
but does not exceed the earlier variants.
Replacing the hinge-style type term with an entropy formulation changes
the smoothness of the optimisation without changing the information
available to the objective,
and similarly produces no meaningful gain.
Combining the two changes does not improve the result.

The narrow performance range across all ten configurations is therefore
the main evidence from the SA-DMoN investigation.
It does not establish that modularity objectives are unsuitable for pose
grouping in general -
a graph whose edges encode learned identity affinity rather than
spatial proximity could behave differently.
It does show that, for the keypoint $k$NN graph used here,
increasingly elaborate modifications to the modularity objective do not
recover the performance of the simpler contrastive-embedding baseline.

Together with the learned-head experiments,
this result substantially narrows the useful design space.
The earlier experiments show that DEC,
slot attention,
and pairwise graph partitioning do not reliably improve grouping once
the representation transfers to COCO.
SA-DMoN then shows that introducing pose structure through a
modularity-based clustering objective does not resolve the problem either.
These results motivate moving the structural information into the
representation-learning stage itself.
PG-GAT follows this direction by retaining the simple $k$-means read-out
and instead changing how information is exchanged between keypoints
while their embeddings are being formed.

\subsection{Why the three PG-GAT modifications are complementary}
\label{sec_disc_embedding_pggat}

The PG-GAT isolation ablation of
Section~\ref{sec_res_pggat_ablations} tests each of the three
skeleton-aware modifications independently before combining them in the
full model.
The three individual configurations are tightly grouped on synthetic
validation at $0.8825$,
$0.8825$, and $0.8834$,
compared with $0.8783$ for the standard GAT baseline.
Enabling all three modifications increases PGA to $0.8896$,
which is $0.0062$ above the strongest individual configuration
and $0.0113$ above the baseline.
The individual differences are small enough that they should not be
interpreted as a ranking of the three mechanisms.
The more useful result is that no single modification accounts for the
performance of the full model.
Their combination produces a larger improvement than any modification
in isolation,
consistent with the three mechanisms providing complementary rather
than redundant information.

This interpretation follows from the different roles of the three
modifications.
The type-pair attention bias introduces learnable additive terms for
the four edge categories used by PG-GAT:
same-type,
skeletal-neighbour,
same-limb,
and cross-body.
In the standard GATv2 formulation,
attention is derived from the features of the incident nodes and does
not explicitly distinguish these structural relationships.
The type-pair term makes the edge category available directly to the
attention calculation,
allowing otherwise similar node pairs to receive different attention
biases according to their relationship in the pose graph.

The skeleton-aware position encoding provides a complementary source of
information.
It exposes L2 distance,
skeleton-hop distance,
and the same-type indicator directly to the attention mechanism.
Although absolute position is already present in the node representation,
a standard GATv2 layer must infer relative geometric relationships
through the interaction between the two node features.
Providing these relationships explicitly reduces the burden on the
attention layer and,
in the case of skeleton-hop distance,
introduces information that cannot be recovered from image coordinates
alone.

The repulsion mechanism addresses a narrower problem.
One of the four attention heads is restricted to same-type edges,
concentrating part of the model capacity on comparisons such as
left-wrist against left-wrist or right-knee against right-knee.
These are particularly informative for person separation because two
detections of the same joint type cannot belong to the same person.
The dedicated head therefore gives the network a separate pathway
in which the contrastive objective can act on the subset of
relationships most directly associated with cross-person discrimination.

The isolation results are consistent with these three mechanisms being
complementary.
Type-pair attention contributes categorical edge structure,
position encoding contributes relative geometric and topological
information,
and the repulsion head reserves capacity for same-type discrimination.
Their individual gains are modest,
but the combined configuration performs better than any of them alone.
Given that these are single-run measurements and the individual
differences are close to the observed run-to-run variation,
the ablation does not establish the precise interaction between their
gradients.
It does, however, support retaining all three modifications in the
architecture that is subsequently taken into the Bayesian sweep.

The architecture sweep then produces a much larger improvement than any
of the isolation changes.
Starting from the full two-layer PG-GAT configuration at $0.8896$,
the sweep over network depth,
attention-head count,
hidden dimension,
output dimension,
joint-type embedding dimension,
and dropout selects a three-layer configuration with a $256$-dimensional
output and reaches $0.9634$ synthetic-validation PGA.
This is an increase of $0.0738$ over the matched-architecture ablation
configuration.
By comparison, the subsequent loss sweep does not find a configuration
that meaningfully exceeds the inherited contrastive settings -
its best result of $0.9617$ is slightly below the architecture-sweep
winner.

The relative scale of these effects gives a more complete interpretation
of the PG-GAT result.
The skeleton-aware modifications provide the inductive biases that
distinguish PG-GAT from the standard GATv2 baseline,
but their benefit depends strongly on the architecture in which they
operate.
The two-layer isolation experiment establishes that the modifications
provide a useful signal at matched compute,
while the architecture sweep shows that considerably more of that signal
can be exploited with an additional layer and a larger embedding space.
In contrast, varying the contrastive margin,
number of repulsion heads,
learning rate,
and weight decay does not produce a comparable gain.

The main lever is therefore not a more elaborate grouping objective
or a finely tuned loss
but the combination of task-specific structure and sufficient
representation capacity.
PG-GAT introduces information that the standard GATv2 formulation does
not expose explicitly,
and the architecture sweep identifies a network size that makes better
use of that information.
The loss sweep, in turn, provides evidence that the improvement is not
primarily the result of a favourable choice of contrastive
hyperparameters.
This distinction is important when interpreting the headline model -
its performance is not attributable to a single architectural
modification or tuned loss value
but to the combination of skeleton-aware message passing and an
architecture sized to exploit it.

\subsection{Synthesis}
\label{sec_disc_embedding_synthesis}

Reading Sections~\ref{sec_disc_embedding_heads}--\ref{sec_disc_embedding_pggat}
together gives a consistent picture of where the main performance
constraint lies.
The learned grouping heads fail for different reasons,
yet none improves on the simple $k$-means baseline after transfer to
COCO.
The SA-DMoN experiments reach the same conclusion from a graph-structural
direction -
changing the null model,
adjacency,
spectral weight,
and type-loss formulation does not overcome the mismatch between
modularity and the spatial $k$NN graph.
In contrast, modifying the embedding network itself produces a measurable
improvement at matched architecture and a much larger gain once the
architecture is sized by the Bayesian sweep.

The evidence therefore identifies the embedding network as the stronger
lever for pose-grouping performance in the setting studied here.
This does not imply that the grouping algorithm is irrelevant
or that a different grouping formulation could never improve the result.
Rather, within the range of grouping methods evaluated in this work,
additional complexity at the read-out stage contributes less than
improving the representation that the read-out receives.
The effectiveness of $k$-means in the final pipeline follows from this
result -
once the embedding separates person identities sufficiently well -
a simple partitioning method is adequate and avoids introducing another
learned component whose behaviour must transfer across domains.

This provides the experimental answer to \emph{RQ1: the lever}.
For the pose-grouping formulation and data considered in this work,
the main opportunity for improvement lies in the construction of the
per-keypoint embedding rather than in replacing its downstream grouping
algorithm.
This finding also motivates the later experiments in the chapter,
which retain $k$-means and investigate how the representation behaves
under domain transfer,
alternative geometries,
and detector-independent evaluation.

The PG-GAT ablation and architecture sweep provide the corresponding
answer to \emph{RQ2: inductive biases}.
Type-pair attention,
skeleton-aware position encoding,
and same-type repulsion each produce a small positive change when
introduced independently at the matched baseline architecture.
Their combination gives the strongest isolation result,
indicating that the three sources of pose structure are complementary
at the resolution of the experiment.
The architecture sweep then shows that the benefit of these inductive
biases depends strongly on the capacity of the network that uses them,
selecting a three-layer,
$256$-dimensional configuration that substantially outperforms the
two-layer ablation model.

Taken together, RQ1 and RQ2 establish the design basis for the
remainder of the discussion.
The first identifies \emph{where} additional modelling effort is most
productive -
the embedding rather than the grouping head.
The second identifies \emph{how} that embedding can be improved -
by exposing pose-specific categorical,
geometric,
and same-type relationships directly to the attention mechanism
and providing sufficient representation capacity to use them.
The subsequent COCO fine-tuning and end-to-end experiments test whether
these gains remain useful once the model leaves the synthetic
distribution on which these design decisions were developed.
